% Generated by src/analysis/code_refs.py ; do not edit.
{\small\setlength{\tabcolsep}{4pt}
\begin{longtable}{@{}l l r p{0.42\textwidth} l@{}}
\caption{Code map: where each mechanism described in the paper is implemented. Line numbers refer to the listings in this appendix.}\label{tab:codemap}\\
\toprule
Segment & Listing & Lines & Description & Section \\
\midrule
\endfirsthead
\toprule
Segment & Listing & Lines & Description & Section \\
\midrule
\endhead
\bottomrule
\endfoot
\multicolumn{5}{@{}l}{\textit{Driver}} \\
\texttt{HeaderBackend} & \ref{lst:header} (\texttt{\detokenize{aio_bench.h}}) & 25-37 & Backend interface: init, issue, wait, fini, compute hints & \S\ref{sec:impl} \\
\texttt{ThreadInit} & \ref{lst:driver} (\texttt{\detokenize{aio_bench.c}}) & 104 & MPI initialized with \texttt{MPI\_THREAD\_MULTIPLE} for every backend & \S\ref{sec:impl} \\
\texttt{TimedMax} & \ref{lst:driver} (\texttt{\detokenize{aio_bench.c}}) & 53-63 & Concurrent kernel run timed as the maximum over ranks & \S\ref{sec:surrogates} \\
\texttt{Calibrate} & \ref{lst:driver} (\texttt{\detokenize{aio_bench.c}}) & 71-89 & Calibration of fixed work to the target $t_{comp}$ & \S\ref{sec:surrogates} \\
\texttt{Baseline} & \ref{lst:driver} (\texttt{\detokenize{aio_bench.c}}) & 148-155 & Calibration call and standalone baseline $t_{comp}^{alone}$ & \S\ref{sec:surrogates} \\
\texttt{Loop} & \ref{lst:driver} (\texttt{\detokenize{aio_bench.c}}) & 158-187 & Timed epoch loop: compute, then issue (sync or async) I/O & \S\ref{sec:workload} \\
\texttt{Reduce} & \ref{lst:driver} (\texttt{\detokenize{aio_bench.c}}) & 189-197 & Per-step maxima over ranks and end-to-end time & \S\ref{sec:metrics} \\
\multicolumn{5}{@{}l}{\textit{Compute surrogates}} \\
\texttt{Sleep} & \ref{lst:kernels} (\texttt{\detokenize{kernels.c}}) & 189-192 & \texttt{sleep}: \texttt{nanosleep} for $t_{comp}$ & \S\ref{sec:surrogates} \\
\texttt{Spin} & \ref{lst:kernels} (\texttt{\detokenize{kernels.c}}) & 25-32 & spin: register-only xorshift loop & \S\ref{sec:surrogates} \\
\texttt{Dgemm} & \ref{lst:kernels} (\texttt{\detokenize{kernels.c}}) & 38-44 & DGEMM: $256^2$ OpenBLAS calls (L2-resident) & \S\ref{sec:surrogates} \\
\texttt{Stream} & \ref{lst:kernels} (\texttt{\detokenize{kernels.c}}) & 50-58 & STREAM triad over $3\times32$\,MiB arrays & \S\ref{sec:surrogates} \\
\texttt{Cart} & \ref{lst:kernels} (\texttt{\detokenize{kernels.c}}) & 64-71 & 3-D periodic Cartesian decomposition & \S\ref{sec:surrogates} \\
\texttt{Exchange} & \ref{lst:kernels} (\texttt{\detokenize{kernels.c}}) & 77-85 & Six-neighbour non-blocking face exchange & \S\ref{sec:surrogates} \\
\texttt{Halo} & \ref{lst:kernels} (\texttt{\detokenize{kernels.c}}) & 87-96 & halo: 512\,KiB face exchange plus allreduce & \S\ref{sec:surrogates} \\
\texttt{HeatHalo} & \ref{lst:kernels} (\texttt{\detokenize{kernels.c}}) & 104-128 & heat: pack, exchange, and unpack ghost faces & \S\ref{sec:surrogates} \\
\texttt{Heat} & \ref{lst:kernels} (\texttt{\detokenize{kernels.c}}) & 130-145 & heat: 7-point Jacobi sweep on $128^3$ cells & \S\ref{sec:surrogates} \\
\texttt{GpuKernel} & \ref{lst:gpu} (\texttt{\detokenize{gpu_kernels.cu}}) & 16-24 & GPU FMA kernel & \S\ref{sec:surrogates} \\
\texttt{GpuSetup} & \ref{lst:gpu} (\texttt{\detokenize{gpu_kernels.cu}}) & 32-40 & Host wait mode: spin vs.\ blocking synchronization & \S\ref{sec:rq4} \\
\texttt{GpuRun} & \ref{lst:gpu} (\texttt{\detokenize{gpu_kernels.cu}}) & 43-50 & GPU work units and stream synchronization & \S\ref{sec:surrogates} \\
\multicolumn{5}{@{}l}{\textit{I/O backends}} \\
\texttt{PosixInit} & \ref{lst:posix} (\texttt{\detokenize{io_posix.c}}) & 31-42 & POSIX AIO with one glibc helper thread & \S\ref{sec:impl} \\
\texttt{PosixWait} & \ref{lst:posix} (\texttt{\detokenize{io_posix.c}}) & 44-57 & Completion of the in-flight \texttt{aio\_write} & \S\ref{sec:impl} \\
\texttt{PosixIssue} & \ref{lst:posix} (\texttt{\detokenize{io_posix.c}}) & 59-90 & Blocking \texttt{pwrite} vs.\ staged \texttt{aio\_write} & \S\ref{sec:impl} \\
\texttt{MpiioIssue} & \ref{lst:mpiio} (\texttt{\detokenize{io_mpiio.c}}) & 60-77 & Collective \texttt{MPI\_File\_write\_at\_all} vs.\ \texttt{iwrite\_at\_all} & \S\ref{sec:impl} \\
\texttt{HdfInit} & \ref{lst:hdf5} (\texttt{\detokenize{io_hdf5.c}}) & 19-42 & MPI-IO VFD, collective transfer, event set, VOL check & \S\ref{sec:impl} \\
\texttt{HdfWait} & \ref{lst:hdf5} (\texttt{\detokenize{io_hdf5.c}}) & 44-51 & \texttt{H5ESwait} on the event set & \S\ref{sec:impl} \\
\texttt{HdfIssue} & \ref{lst:hdf5} (\texttt{\detokenize{io_hdf5.c}}) & 53-81 & Per-step dataset create/write/close via the \texttt{\_async} API & \S\ref{sec:impl} \\
\texttt{BpInit} & \ref{lst:adios2} (\texttt{\detokenize{io_adios2.c}}) & 27-61 & BP5 setup: \texttt{AsyncWrite} mode and aggregation parameters & \S\ref{sec:impl} \\
\texttt{BpIssue} & \ref{lst:adios2} (\texttt{\detokenize{io_adios2.c}}) & 63-73 & One ADIOS2 step per timestep (deferred \texttt{Put}) & \S\ref{sec:impl} \\
\texttt{BpHints} & \ref{lst:adios2} (\texttt{\detokenize{io_adios2.c}}) & 85-97 & Computation-block hints for Guided mode & \S\ref{sec:impl} \\
\texttt{PdcIssue} & \ref{lst:pdc} (\texttt{\detokenize{io_pdc.c}}) & 48-62 & PDC region transfer (implemented, not evaluated) & \S\ref{sec:impl} \\
\multicolumn{5}{@{}l}{\textit{Experiment scripts}} \\
\texttt{Binding} & \ref{lst:runone} (\texttt{\detokenize{run_one.sh}}) & 14-21 & Rank binding: fully subscribed, SMT sibling, core reserved & \S\ref{sec:machine} \\
\texttt{AsyncEnv} & \ref{lst:runone} (\texttt{\detokenize{run_one.sh}}) & 27-36 & Async VOL selection, \texttt{EXE\_DCLOSE}, collective HDF5 & \S\ref{sec:impl} \\
\texttt{Cgroup} & \ref{lst:runone} (\texttt{\detokenize{run_one.sh}}) & 40-44 & \texttt{sync}, cgroup memory limit, and launch & \S\ref{sec:machine} \\
\texttt{MatrixEnv} & \ref{lst:runmatrix} (\texttt{\detokenize{run_matrix.py}}) & 27-32 & Library settings shared by all runs & \S\ref{sec:impl} \\
\texttt{MemMax} & \ref{lst:runmatrix} (\texttt{\detokenize{run_matrix.py}}) & 42-45 & Per-run cgroup limit: footprint + 24\,GiB & \S\ref{sec:machine} \\
\texttt{IoBaseline} & \ref{lst:runmatrix} (\texttt{\detokenize{run_matrix.py}}) & 101-106 & Per-phase I/O-only baseline that sizes $t_{comp}$ & \S\ref{sec:workload} \\
\texttt{TioSync} & \ref{lst:runmatrix} (\texttt{\detokenize{run_matrix.py}}) & 80-98 & Synchronous I/O time per step from the baseline runs & \S\ref{sec:workload} \\
\texttt{RunOne} & \ref{lst:runmatrix} (\texttt{\detokenize{run_matrix.py}}) & 60-77 & Resumable execution of one configuration & \S\ref{sec:workload} \\
\texttt{PhaseMain} & \ref{lst:runmatrix} (\texttt{\detokenize{run_matrix.py}}) & 125-132 & Main experiment: all CPU surrogates, one run each & \S\ref{sec:rq1} \\
\texttt{PhaseSize} & \ref{lst:runmatrix} (\texttt{\detokenize{run_matrix.py}}) & 135-142 & Size sweep, 0.75 to 48\,GiB per step & \S\ref{sec:workload} \\
\texttt{PhaseSizeLarge} & \ref{lst:runmatrix} (\texttt{\detokenize{run_matrix.py}}) & 189-196 & 48\,GiB per step for POSIX and ADIOS2 & \S\ref{sec:workload} \\
\texttt{PhaseSweep} & \ref{lst:runmatrix} (\texttt{\detokenize{run_matrix.py}}) & 155-163 & $t_{comp}/t_{io}$ sweep & \S\ref{sec:rq2} \\
\texttt{PhaseBind} & \ref{lst:runmatrix} (\texttt{\detokenize{run_matrix.py}}) & 166-177 & Core-allocation experiment & \S\ref{sec:rq4} \\
\texttt{PhaseGpu} & \ref{lst:runmatrix} (\texttt{\detokenize{run_matrix.py}}) & 180-186 & GPU experiment (4 ranks) & \S\ref{sec:rq4} \\
\texttt{RunPdc} & \ref{lst:runmatrix} (\texttt{\detokenize{run_matrix.py}}) & 212-244 & One PDC run, killed as a group after 30 minutes & \S\ref{sec:impl} \\
\texttt{PhasePdc} & \ref{lst:runmatrix} (\texttt{\detokenize{run_matrix.py}}) & 250-268 & PDC main experiment and SMT-sibling placement & \S\ref{sec:rq4} \\
\texttt{PhasePdcSize} & \ref{lst:runmatrix} (\texttt{\detokenize{run_matrix.py}}) & 288-297 & PDC size sweep, 0.73 and 2.9\,GiB per step & \S\ref{sec:size} \\
\texttt{PhasePdcSweep} & \ref{lst:runmatrix} (\texttt{\detokenize{run_matrix.py}}) & 300-308 & PDC $t_{comp}/t_{io}$ sweep, 0.25 to 1 & \S\ref{sec:sweep} \\
\texttt{PhasePdcGpu} & \ref{lst:runmatrix} (\texttt{\detokenize{run_matrix.py}}) & 278-285 & PDC GPU experiment (1 server, 4 clients) & \S\ref{sec:gpu} \\
\multicolumn{5}{@{}l}{\textit{Analysis}} \\
\texttt{Metrics} & \ref{lst:plotsmeasured} (\texttt{\detokenize{plots_measured.py}}) & 88-113 & Speedup, compute slowdown, and exposed I/O & \S\ref{sec:metrics} \\
\texttt{Excess} & \ref{lst:plotsmeasured} (\texttt{\detokenize{plots_measured.py}}) & 222-245 & Run time relative to the sleep prediction & \S\ref{sec:metrics} \\
\multicolumn{5}{@{}l}{\textit{Build and patches}} \\
\texttt{BuildHdf} & \ref{lst:builddeps} (\texttt{\detokenize{build_deps.sh}}) & 22-30 & Parallel, thread-safe HDF5 build & \S\ref{sec:machine} \\
\texttt{BuildVol} & \ref{lst:builddeps} (\texttt{\detokenize{build_deps.sh}}) & 32-41 & Async VOL build with the wait patch & \S\ref{sec:impl} \\
\texttt{BuildAdios} & \ref{lst:builddeps} (\texttt{\detokenize{build_deps.sh}}) & 43-56 & ADIOS2 build options & \S\ref{sec:machine} \\
\texttt{VolPatch} & \ref{lst:volasyncpatch} (\texttt{\detokenize{vol-async-wait.patch}}) & 1-22 & Async VOL wait back-off: 1\,s and 100\,ms $\to$ 1\,ms & \S\ref{sec:impl} \\
\texttt{AdiosPatch} & \ref{lst:adiospatch} (\texttt{\detokenize{adios2-fsync.patch}}) & 1-56 & Opt-in per-step \texttt{fdatasync} for BP5 (inactive) & \S\ref{sec:machine} \\
\end{longtable}}
